problem:  
Let \((X, d, \mathfrak{m})\) be an \(\mathrm{RCD}(K,N)\) metric measure space with \(K>0\) and \(1<N<\infty\).  
Denote by \(C_{N,K}^{\mathrm{sphere}}\) the sharp Sobolev constant attained on the canonical \(N\)-sphere of Ricci curvature \(K\), so that for the round sphere \((S^N, g_{\mathrm{round}}, \mathrm{vol}_{S^N})\) one has  
\[
\|f\|_{L^{\frac{2N}{N-2}}(\mathrm{vol}_{S^N})}^2 
\le C_{N,K}^{\mathrm{sphere}}\left(
   \|\nabla f\|_{L^2(\mathrm{vol}_{S^N})}^2 + 
   \frac{N-2}{N} K \|f\|_{L^2(\mathrm{vol}_{S^N})}^2
\right).
\]  

It is known that such inequalities, when formulated for \(\mathrm{RCD}(K,N)\) spaces, imply the spectral gap and isoperimetric inequalities from which *rigidity for equality cases* on spheres follows in several instances (e.g., Cavalletti–Mondino, Ketterer). Yet all existing rigidity results treat inequalities derived from linear or isoperimetric functionals rather than the *nonlinear Sobolev embedding constant itself*.  

**Open Problem:**  
Investigate the *quantitative stability and extremal structure* of the Sobolev inequality in the \(\mathrm{RCD}(K,N)\) setting:  
Can one prove the existence of a universal, nontrivial function \(\Phi(\varepsilon)>0\) with \(\Phi(\varepsilon)\to 0\) as \(\varepsilon\to 0\) such that  
for every \(\mathrm{RCD}(K,N)\) space \((X,d,\mathfrak{m})\) satisfying  
\[
\sup_{f\ne 0}
\frac{\|f\|_{L^{\frac{2N}{N-2}}(\mathfrak{m})}^2}
     {\|Df\|_{L^2(\mathfrak{m})}^2 + \frac{N-2}{N}K\|f\|_{L^2(\mathfrak{m})}^2}
\ge (1-\varepsilon)\, C_{N,K}^{\mathrm{sphere}},
\]
the space is \(\Phi(\varepsilon)\)-close in the measured Gromov–Hausdorff topology to a spherical suspension over an \(\mathrm{RCD}(N-2,N-1)\) space?  

Equivalently:  
Is there a *quantitative rigidity* principle for the sharp Sobolev constant in the \(\mathrm{RCD}(K,N)\) class, describing how closeness to optimal constant enforces geometric closeness to the round sphere?  

Why is it a "good" problem:  
This problem refines classical Obata and Lévy–Gromov rigidity from binary equality-versus-none dichotomies into *quantitative metric analysis* within the nonlinear functional setting of Sobolev constants on nonsmooth spaces.  
It is deeply open: while equality cases in certain sharp inequalities are well understood, *stability near equality* for Sobolev-type embeddings in \(\mathrm{RCD}(K,N)\) spaces remains completely unexplored. The question fuses **geometric analysis** (Sobolev constants and extremizers), **optimal transport theory** (synthetic Ricci curvature), and **metric geometry** (measured Gromov–Hausdorff convergence).  
A solution would create a powerful bridge between analysis and geometry, providing a new rigidity–stability paradigm for nonlinear inequalities in synthetic curvature settings. The formulation is conceptually simple yet demands deep, multi-field insight—an archetype of a profound and forward-looking research problem.